\documentclass[a4paper,11pt]{article}
\usepackage[italian]{babel}
\usepackage[T1]{fontenc}
\usepackage[utf8]{inputenc}
\usepackage{lmodern}
\usepackage{geometry}
\usepackage{xcolor}
\usepackage{hyperref}
\usepackage{tabularx}
\usepackage{array}
\usepackage{booktabs}
\usepackage{enumitem}
\usepackage{titlesec}
\usepackage{tcolorbox}
\usepackage{fancyhdr}
\usepackage{multicol}
\usepackage{amsmath}

\geometry{margin=2.2cm}
\setlength{\parskip}{0.6em}
\setlength{\parindent}{0pt}
\setlength{\headheight}{14pt}

% --- Colori ---
\definecolor{MainColor}{HTML}{1F4E79}
\definecolor{AccentColor}{HTML}{2E86AB}
\definecolor{SoftBlue}{HTML}{EAF4FB}
\definecolor{SoftGray}{HTML}{F5F7FA}
\definecolor{SoftGreen}{HTML}{EAF8F0}
\definecolor{DarkText}{HTML}{1E1E1E}

% --- Hyperref ---
\hypersetup{
    colorlinks=true,
    linkcolor=MainColor,
    urlcolor=AccentColor,
    pdftitle={Guida moderna a SSH},
    pdfauthor={}
}

% --- Titoli ---
\titleformat{\section}
  {\Large\bfseries\color{MainColor}}
  {\thesection}{0.8em}{}

\titleformat{\subsection}
  {\large\bfseries\color{AccentColor}}
  {\thesubsection}{0.8em}{}

% --- Header/Footer ---
\pagestyle{fancy}
\fancyhf{}
\fancyhead[L]{\textcolor{MainColor}{Guida a SSH}}
\fancyhead[R]{\textcolor{AccentColor}{Accesso remoto e tunnel}}
\fancyfoot[C]{\thepage}

% --- Box ---
\newtcolorbox{infobox}{
    colback=SoftBlue,
    colframe=AccentColor,
    boxrule=0.8pt,
    arc=3mm,
    left=2mm,right=2mm,top=2mm,bottom=2mm
}

\newtcolorbox{tipbox}{
    colback=SoftGreen,
    colframe=green!50!black,
    boxrule=0.8pt,
    arc=3mm,
    left=2mm,right=2mm,top=2mm,bottom=2mm
}

\newtcolorbox{codebox}{
    colback=SoftGray,
    colframe=black!20,
    boxrule=0.5pt,
    arc=2mm,
    left=2mm,right=2mm,top=2mm,bottom=2mm
}

\begin{document}

% =========================
% Titolo
% =========================
\vspace*{1.5cm}

{\Huge\bfseries\color{MainColor} Guida a \texttt{SSH}} \\[0.4cm]
{\Large\color{AccentColor} Connessioni sicure, chiavi, config e port forwarding}

\vspace{1cm}

\begin{infobox}
\textbf{SSH (Secure Shell)} e il modo standard per collegarti a macchine remote in modo cifrato: login shell, esecuzione comandi, copia file e tunnel di rete.
\end{infobox}

\vspace{0.5cm}

\begin{codebox}
\textbf{Quick reference}

\renewcommand{\arraystretch}{1.2}
\begin{tabularx}{\textwidth}{>{\ttfamily}l X}
\toprule
Comando & Descrizione \\
\midrule
ssh user@host & Connessione base \\
ssh -p 2222 user@host & Connessione su porta diversa \\
ssh -i \textasciitilde/.ssh/id\_ed25519 user@host & Specifica chiave \\
ssh -v user@host & Debug connessione (verbose) \\
ssh-keygen -t ed25519 & Genera coppia di chiavi \\
ssh-add \textasciitilde/.ssh/id\_ed25519 & Carica chiave nell'agent \\
scp file user@host:/path & Copia file su remoto \\
ssh -L 8080:localhost:80 user@host & Tunnel locale (port forward) \\
ssh -J bastion user@target & Jump host / bastion \\
\bottomrule
\end{tabularx}
\end{codebox}

\vspace{0.6cm}
\tableofcontents
\newpage

% =========================
\section{Cos'e SSH}
\texttt{SSH} e un protocollo che ti permette di stabilire una connessione cifrata tra un client (il tuo computer) e un server remoto.

In pratica lo userai per:
\begin{itemize}[leftmargin=1.2cm]
    \item entrare su un server e lavorare da shell;
    \item eseguire comandi remoti;
    \item trasferire file (\texttt{scp}, \texttt{sftp}, \texttt{rsync -e ssh});
    \item creare tunnel di rete (port forwarding) per accedere a servizi non esposti.
\end{itemize}

\begin{tipbox}
\textbf{Idea pratica:} se impari bene SSH, diventa la "colla" tra sviluppo locale, server e strumenti (Git, deploy, DB, dashboard interne).
\end{tipbox}

% =========================
\section{Concetti fondamentali}
\begin{infobox}
\textbf{Client / Server} \\
Il client e il comando \texttt{ssh} sul tuo computer. Il server e il demone \texttt{sshd} sulla macchina remota.

\medskip
\textbf{Porta} \\
Di default SSH usa la porta \texttt{22}, ma spesso viene cambiata (es. \texttt{2222}).

\medskip
\textbf{Chiave pubblica / privata} \\
Autenticazione moderna: la chiave privata resta sul client; la pubblica viene messa sul server in \texttt{\textasciitilde/.ssh/authorized\_keys}.

\medskip
\textbf{known\_hosts} \\
Il file in cui il client salva le impronte dei server visitati per evitare attacchi man-in-the-middle.
\end{infobox}

Struttura mentale:
\[
\text{Client} \rightarrow \text{Server} \rightarrow \text{Autenticazione (password o chiavi)}
\]

% =========================
\section{Primo utilizzo}
\subsection{Connessione base}
\begin{codebox}
\texttt{ssh user@host}
\end{codebox}

Dove:
\begin{itemize}[leftmargin=1.2cm]
    \item \texttt{user} e l'utente sul server (es. \texttt{ubuntu}, \texttt{root}, \texttt{deploy});
    \item \texttt{host} e un nome DNS o un IP.
\end{itemize}

\subsection{Porta diversa}
\begin{codebox}
\texttt{ssh -p 2222 user@host}
\end{codebox}

\subsection{Usare una chiave specifica}
\begin{codebox}
\texttt{ssh -i \textasciitilde/.ssh/id\_ed25519 user@host}
\end{codebox}

\begin{tipbox}
\textbf{Consiglio:} se stai usando chiavi, punta a non passare \texttt{-i} ogni volta: metti la configurazione in \texttt{\textasciitilde/.ssh/config}.
\end{tipbox}

% =========================
\section{Chiavi SSH (setup consigliato)}
\subsection{Generare una chiave moderna (ed25519)}
\begin{codebox}
\texttt{ssh-keygen -t ed25519 -C "nome@azienda.com"}
\end{codebox}

Di default troverai:
\begin{itemize}[leftmargin=1.2cm]
    \item chiave privata: \texttt{\textasciitilde/.ssh/id\_ed25519}
    \item chiave pubblica: \texttt{\textasciitilde/.ssh/id\_ed25519.pub}
\end{itemize}

\begin{infobox}
\textbf{Passphrase:} mettila. E una protezione importante se la chiave privata viene copiata o rubata.
\end{infobox}

\subsection{Caricare la chiave nell'agent}
\begin{codebox}
\texttt{eval "\$(ssh-agent -s)"} \\
\texttt{ssh-add \textasciitilde/.ssh/id\_ed25519}
\end{codebox}

Su macOS spesso e utile anche salvare la chiave nel portachiavi (opzionale):
\begin{codebox}
\texttt{ssh-add --apple-use-keychain \textasciitilde/.ssh/id\_ed25519}
\end{codebox}

% =========================
\section{Aggiungere la chiave al server}
\subsection{Metodo 1: ssh-copy-id (se disponibile)}
\begin{codebox}
\texttt{ssh-copy-id -i \textasciitilde/.ssh/id\_ed25519.pub user@host}
\end{codebox}

\subsection{Metodo 2: senza ssh-copy-id (compatibile ovunque)}
\begin{codebox}
\texttt{cat \textasciitilde/.ssh/id\_ed25519.pub | ssh user@host "mkdir -p \textasciitilde/.ssh \&\& cat >> \textasciitilde/.ssh/authorized\_keys"}
\end{codebox}

\begin{tipbox}
\textbf{Check veloce:} dopo aver aggiunto la chiave, riprova \texttt{ssh user@host}. Se tutto e ok, non ti chiedera la password (ma potrebbe chiedere la passphrase della chiave).
\end{tipbox}

% =========================
\section{Configurare \texttt{\textasciitilde/.ssh/config}}
Con \texttt{\textasciitilde/.ssh/config} puoi creare alias e standardizzare porte, user e chiavi.

Esempio:
\begin{codebox}
\texttt{Host myserver} \\
\texttt{  HostName 203.0.113.10} \\
\texttt{  User ubuntu} \\
\texttt{  Port 2222} \\
\texttt{  IdentityFile \textasciitilde/.ssh/id\_ed25519} \\
\texttt{  IdentitiesOnly yes}
\end{codebox}

Da quel momento:
\begin{codebox}
\texttt{ssh myserver}
\end{codebox}

\begin{infobox}
\textbf{IdentitiesOnly yes} evita tentativi con troppe chiavi (errore tipico: \texttt{Too many authentication failures}).
\end{infobox}

% =========================
\section{Copia file: scp e rsync}
\subsection{scp (semplice)}
\begin{codebox}
\texttt{scp file.txt user@host:/tmp/} \\
\texttt{scp -r cartella/ user@host:/var/www/}
\end{codebox}

\subsection{rsync (piu efficiente su cartelle grandi)}
\begin{codebox}
\texttt{rsync -av --progress -e ssh cartella/ user@host:/var/www/}
\end{codebox}

% =========================
\section{Port forwarding (tunnel)}
Il port forwarding serve per "portare" un servizio remoto in locale (o viceversa) senza esporlo pubblicamente.

\subsection{Tunnel locale (\texttt{-L})}
Esempio: hai un servizio sul server su \texttt{localhost:80} e vuoi vederlo in locale su \texttt{localhost:8080}.
\begin{codebox}
\texttt{ssh -L 8080:localhost:80 user@host}
\end{codebox}

\subsection{Tunnel remoto (\texttt{-R})}
Esempio: vuoi esporre una porta locale (es. \texttt{3000}) sul server remoto.
\begin{codebox}
\texttt{ssh -R 3000:localhost:3000 user@host}
\end{codebox}

\subsection{SOCKS proxy (\texttt{-D})}
Crea un proxy SOCKS locale (utile per navigare tramite il server):
\begin{codebox}
\texttt{ssh -D 1080 user@host}
\end{codebox}

\begin{tipbox}
\textbf{Consiglio:} aggiungi \texttt{-N} quando ti serve solo il tunnel (senza aprire shell).
\end{tipbox}

% =========================
\section{Jump host (bastion)}
Quando il server target non e direttamente raggiungibile, passi da un bastion.

\subsection{Usare \texttt{-J}}
\begin{codebox}
\texttt{ssh -J user@bastion user@target}
\end{codebox}

\subsection{Configurazione con ProxyJump}
\begin{codebox}
\texttt{Host bastion} \\
\texttt{  HostName 198.51.100.10} \\
\texttt{  User ubuntu} \\
\texttt{Host target} \\
\texttt{  HostName 10.0.1.25} \\
\texttt{  User ubuntu} \\
\texttt{  ProxyJump bastion}
\end{codebox}

% =========================
\section{Troubleshooting veloce}
\subsection{Debug con verbose}
\begin{codebox}
\texttt{ssh -v user@host} \\
\texttt{ssh -vvv user@host}
\end{codebox}

\subsection{Errori tipici}
\begin{infobox}
\textbf{Permission denied (publickey)} \\
La chiave non e sul server, oppure stai usando la chiave sbagliata.

\medskip
\textbf{Host key verification failed} \\
Il server e cambiato (o la chiave in \texttt{known\_hosts} non coincide). Verifica prima di rimuovere l'entry.

\medskip
\textbf{Connection refused / timed out} \\
Porta sbagliata, firewall, o \texttt{sshd} non in ascolto.
\end{infobox}

\subsection{Permessi della cartella \texttt{.ssh}}
Su molti server, permessi troppo aperti bloccano l'autenticazione a chiave.
\begin{codebox}
\texttt{chmod 700 \textasciitilde/.ssh} \\
\texttt{chmod 600 \textasciitilde/.ssh/authorized\_keys}
\end{codebox}

% =========================
\section{Cheat sheet veloce}
\begin{center}
\renewcommand{\arraystretch}{1.25}
\begin{tabularx}{\textwidth}{>{\ttfamily}l X}
\toprule
Comando & Significato \\
\midrule
ssh user@host & Login remoto \\
ssh -p 2222 user@host & Porta custom \\
ssh -i key user@host & Specifica chiave \\
ssh -v user@host & Verbose (debug) \\
ssh-keygen -t ed25519 & Genera chiave \\
ssh-add key & Carica chiave nell'agent \\
scp file user@host:/path & Copia file \\
rsync -e ssh src/ user@host:dst/ & Sync efficiente \\
ssh -L lport:host:rport user@host & Tunnel locale \\
ssh -R rport:host:lport user@host & Tunnel remoto \\
ssh -D 1080 user@host & SOCKS proxy \\
ssh -J bastion user@target & Jump host \\
\bottomrule
\end{tabularx}
\end{center}

% =========================
\section{Errori comuni dei principianti}
\begin{infobox}
\textbf{1. Riutilizzare la stessa chiave ovunque} \\
Meglio chiavi separate per contesti diversi (lavoro, personale, produzione).

\medskip
\textbf{2. Usare \texttt{root} direttamente} \\
Preferisci un utente normale + \texttt{sudo} (quando serve).

\medskip
\textbf{3. Disabilitare i controlli di sicurezza a caso} \\
Evitare di forzare \texttt{StrictHostKeyChecking no} senza capire cosa implica.
\end{infobox}

% =========================
\section{Conclusione}
\texttt{SSH} e uno strumento semplice all'apparenza, ma potentissimo quando lo usi bene: chiavi + config + tunnel.

\begin{tipbox}
\textbf{Da ricordare:}
\begin{itemize}[leftmargin=1.2cm]
    \item usa chiavi \texttt{ed25519} con passphrase;
    \item centralizza alias e opzioni in \texttt{\textasciitilde/.ssh/config};
    \item usa i tunnel per accedere a servizi interni senza esporli;
    \item se qualcosa non funziona: \texttt{ssh -v} e controlla permessi.
\end{itemize}
\end{tipbox}

\vfill
\begin{center}
\textcolor{AccentColor}{\large Buon lavoro con \texttt{SSH}}
\end{center}

\end{document}
